% A report fragment: requires amsmath, amssymb, amsthm.
\section{The crossed no-push radix construction}

For $k\in\{5,6\}$, let
\[
 p(a,b)=\frac{(a+b)(a+b+1)}2+b,
 \quad I_k(A,B)=k p(A,B),\quad M_k(a,b)=k p(a,b)+1,
\]
\[
 U_k(a,B)=k p(B,a)+2,\qquad E_k(a,b)=k p(a,b)+3.
\]
Numerical code zero is truth and code one is the designated contradiction.
Unused residues are decoded as truth. These definitions retain the exceptional
identity $I_k(0,0)=0$; no injectivity assertion separates this instance from truth.
For nonzero constructor codes, every recursive formula child is smaller than
its parent, and every occurring variable label is smaller than the whole code.
The induction proving this uses $a,b\le p(a,b)$ and $k\ge4$.

The numerical calculus uses the separate diagonal head
$\Gamma_k(B)=I_k(B,U_k(B,B))$ and the inherited separate-head gate.
Its distribution head is the exchanged version
\[
 U_k(d,I_k(a,h))\ \longrightarrow\
       \bigl(U_k(d,a)\longrightarrow U_k(d,h)\bigr).
\]
The other heads and rules are those of the supplied radix-family calculus.
The fixed implicational exchange derivation converts this head into the ordinary
distribution head, and conversely, using only the derived $K,S$ basis and modus
ponens. The resulting proof transformations are primitive recursive.

The inherited uniform constructor-$k$ calibration consequently applies. One
must not replace this argument by a claim that code substitution alone conjugates
the two calculi: the diagonal binder is the whole formula code, and hence changes
when $k$ changes. The uniform calibration instead uses the displayed strict
fresh-label bound and instantiates the proof-producing PA templates separately
at each coefficient. This is an explicitly inherited ordinary mathematical
dependency, not a newly compiled proof-assistant theorem.

\subsection{No-push commands and the all-store compiler}

Write $[a;w]=p(a,w)$ for the numerical stack. Its projections are total, including
$[0;0]=0$ and $\operatorname{hd}(0)=\operatorname{tl}(0)=0$. For a command store
$(t,w,d,h)$, the modified modus ponens command is
\[
 \mathsf{MP}_x(t,w,d,h)=
 \begin{cases}
   (b,w,b,0),&I_k(t,b)=h,\\
   (0,w,b,0),&I_k(t,b)\ne h,
 \end{cases}\qquad b=d+x.
\]
The Shift command retains its sum test: it halts if $t+h=1$ and otherwise returns
$(h,\operatorname{tl}(w),d,\operatorname{hd}(w))$. Thus this construction differs
from the cached-head-only machines as well as from the dummy-push radix machines.

The AX command pushes nine generated entries and the saved old top. Its last
generated entry in pop order is the universal code $U_k(d,h)\ge2$. The other
generated entries and the guarded top have implication residue zero; the
universal entry has residue two. None can equal one, independently of consistency.
The exchanged distribution head changes its formula but not these facts or its
position in the ten-entry protocol.

Set $A_0=\mathsf{AX}_{0,0,0}$ and use the supplied no-push compiler words
\[
 \mathsf{Push}=A_0;S^9;\mathsf{MP}_0,
 \qquad \mathsf{Pop}=A_0;S^{11},
 \qquad \mathsf{Discard}=\mathsf{MP}_0.
\]
They satisfy, uniformly in $k\ge4$,
\[
 \mathsf{Push}(t,w,d,h)=(0,[t;w],0,0),\qquad
 \mathsf{Pop}(t,[a;w],d,h)=(t,w,0,a),
\]
\[
 \mathsf{Discard}(t,w,0,h)=(0,w,0,0).
\]
For Push the final implication comparison fails because residues zero and two
are disjoint. The potentially last Shift in Pop tests $U_k(d,h)+t\ge2$.
Discard is valid because its success and failure outcomes coincide at $d=x=0$.
These are numerical equalities even when zero stack entries collapse.

The axiom selectors and generalization word of the supplied no-push compiler
therefore remain valid. For modus ponens, reserve a zero below the two premises
and execute
\[
 \mathsf{Discard};\mathsf{Push};C(I_k(a,b));\mathsf{Push};
 C(a);\mathsf{Push};S;S;\mathsf{MP}_b;\mathsf{Pop}.
\]
After successful recursive calls, the store before the two Shifts is
$(0,[a,I_k(a,b),0;w],0,0)$. The second Shift may halt early at $a=1$;
otherwise MP returns $(b,[0;w],b,0)$ and Pop returns $(b,w,0,0)$.
Thus induction on a finite derivation shows that its compiled word either
halts or returns its conclusion with the surrounding stack unchanged. Prefixing
by $A_0$ and appending Shift produces a halting word for every incoming store
when the derivation concludes one. The inherited cofinal-word scheduler proof
applies; it never assumes that unsuccessful trials reset their theorem store.

Conversely, attach a derivation to the visible top, cached head, and every live
canonical stack entry. AX, MP, and Shift preserve this annotation. If Shift
halts, $t+h=1$ implies that one of the two annotated entries equals one. The
inherited syntactic interpretation turns this into a PA refutation. Hence the
source halting predicate and the PA-refutation predicate are PA-provably equivalent.

\subsection{Binding to literal machines}

The construction uses the pinned Q AX grammar, the no-push replacement
$\texttt{node163}=(\texttt{node6},\texttt{node54})$, and the multiplier
\[
 \texttt{node61}=
 \begin{cases}
 ((\mathsf{INC}_1,\texttt{node60}),\texttt{node60}),&k=5,\\
 ((\mathsf{INC}_1,\texttt{node60}),(\mathsf{INC}_1,\texttt{node60})),&k=6.
 \end{cases}
\]
There are exactly 61 multiplier occurrences. The source-binding certificate
checks their complete Cantor-pairing prefixes, changes each multiplication body
from four to $k$ increments, and proves that all exterior source edges agree
under the explicit relocation. No exterior edge enters a changed body interior.
The complete AX macro word is checked in a free constructor algebra.

Each raw table has 202 controls. Seven Q-reader folds and the additional
$41\mapsto99$ fold give 194 controls; all 32 local cases are checked in the final
simultaneous quotient. The 23 arithmetic-engine controls retain every baseline
row. Complete blank initialization takes 58,320 literal transitions and reaches
the clean source frame, including all nonzero cells and the blank exteriors.

The constructor checker verifies every primitive source outcome for arbitrary
natural register values; the separate literal checker reconstructs the four
list lemmas from the transition rows and resolves the source in two ways.
These finite certificates support elementary induction arguments for the
source/tape simulation. Positive simulation durations give halting reflection:
if the literal machine first halted at time $T$ but the source had not halted
within $T+1$ steps, their deterministic positive simulations would continue
beyond that first halt, a contradiction. All finite traces and the bounded
iteration are arithmetically representable in PA.

Composing the two source reductions with these source/tape simulations yields
the ordinary mathematical statement
\[
 \mathrm{PA}\vdash
   \neg\operatorname{Halt}(\mathrm{PA194QN}k,\mathrm{blank})
       \ \longleftrightarrow\ \operatorname{Con}(\mathrm{PA}),
 \qquad k=5,6.
\]
The new result is this explicit composition and its literal bindings. None of
its component optimization principles is claimed as new. A successful executable
certificate is not an end-to-end Lean/Rocq verification or a completed iPRA endpoint.
